%! Potential Errors When Performing Tests — Unit 6, Lesson 6.7 — Homework
\documentclass[10pt]{article}
\usepackage{apstats-article}
\usepackage{apstats-boxes}
\newcommand{\TallMath}[1]{$\displaystyle #1\rule[-1.4em]{0pt}{3.2em}$}

\begin{document}

\pageheader{Unit 6, Lesson 6.7}{Homework}
\namedateperiod

\vspace{0.15in}

\begin{enumerate}

  \item A city health department tests $H_0: p = 0.10$ (10\% of adults have high blood
        pressure) vs.\ $H_a: p > 0.10$.
        \begin{enumerate}[label=\alph*.]
          \item Describe a Type I error in context. \writelines{2}
          \item Describe a Type II error in context. \writelines{2}
          \item If the true proportion is 12\%, would power be higher or lower than if
                the true proportion were 20\%? Explain. \writelines{2}
        \end{enumerate}

  \item A new reading intervention is tested at a school. $H_0$: the intervention does not
        improve reading scores; $H_a$: the intervention improves reading scores.
        \begin{enumerate}[label=\alph*.]
          \item From the perspective of a school principal, which error would be more
                costly? Explain. \writelines{3}
          \item The school increases the sample from 60 students to 240 students. How
                does this change the power of the test and Type II error probability? \writelines{2}
        \end{enumerate}

  \item \textbf{True or False.} For each statement, write T or F and explain.
        \begin{enumerate}[label=\alph*.]
          \item ``Reducing $\alpha$ from 0.05 to 0.01 always increases power.''
                \blank{1cm} Explain: \writelines{2}
          \item ``A Type II error can only occur if $H_0$ is rejected.''
                \blank{1cm} Explain: \writelines{2}
        \end{enumerate}

\end{enumerate}

\begin{reflectionbox}
\textbf{Preview:} Lesson 6.8 introduces the two-sample $z$-interval for the difference
of two proportions $p_1 - p_2$ --- when we want to compare two groups rather than test
a single proportion against a benchmark.
\end{reflectionbox}

\end{document}
